\chapter{Von der Idee zum Werkstück}

In diesem Kapitel wird der gesamte Ablauf des CNC-Hot-Wire-Foam-Cutters
"ELMAR" beschrieben -- vom ersten Entwurf einer Kontur bis zum fertig
geschnittenen Werkstück aus Schaumstoff. Die vorangegangenen Kapitel
haben die einzelnen Komponenten des Systems im Detail behandelt: den
Arduino Uno als Steuerzentrale mit der GRBL-Firmware, das CNC Shield V3
mit den Schrittmotortreibern, die Endschalter zur Positionsbestimmung,
den Druckknopf zur Bedienung, das MOSFET-PWM-Schaltmodul zur Ansteuerung
des Heizdrahtes sowie den Universal G-Code Sender (UGS) als
Bediensoftware auf dem angeschlossenen Computer. Dieses Kapitel führt
diese Bausteine zu einem durchgängigen Arbeitsablauf zusammen.

\section{Idee und Entwurf der Kontur}

Am Anfang jedes Werkstücks steht die Idee, welche Form aus dem
Schaumstoffblock herausgeschnitten werden soll. Diese Form wird zunächst
als zweidimensionale Kontur entworfen, beispielsweise als Skizze oder
als digitale Zeichnung. Da der Heizdraht beim Schneiden entlang einer
durchgehenden Bahn geführt wird, muss die Kontur als geschlossener oder
durchgehend befahrbarer Linienzug vorliegen, ohne Sprungstellen, die der
Draht nicht abfahren kann.

Bei der Erstellung der Kontur ist außerdem die Größe des verwendeten
Schaumstoffblocks sowie der Verfahrbereich der Achsen zu
berücksichtigen. Die Kontur darf den durch die Endschalter begrenzten
Arbeitsbereich der Maschine nicht überschreiten, da es sonst beim
späteren Schneidvorgang zu einer Kollision mit den mechanischen
Endanschlägen kommen kann.

\section{Erzeugung der Steuerbefehle (G-Code)}

Aus der entworfenen Kontur werden anschließend die Steuerbefehle für die
Maschine abgeleitet. Diese liegen im G-Code-Format vor, einer
standardisierten Befehlssprache für die Steuerung von CNC-Maschinen. Der
G-Code beschreibt die einzelnen Bewegungen der Achsen in Form von
Koordinaten und Verfahrgeschwindigkeiten, die der Heizdraht beim
Schneidvorgang nacheinander abfährt. Der G-Code kann dabei entweder
manuell geschrieben oder mithilfe einer CAM-Software aus einem
CAD-Modell erzeugt und als Datei (z. B. .nc oder .gcode) exportiert
werden.

Die Umsetzung der G-Code-Befehle in konkrete Schrittmotorsignale
übernimmt die auf dem Arduino Uno installierte GRBL-Firmware. GRBL
interpretiert die im G-Code enthaltenen Bewegungsanweisungen und erzeugt
daraus die STEP- und DIR-Signale, die über das CNC Shield V3 an die
Motortreiber der einzelnen Achsen weitergegeben werden. Auf diese Weise
wird aus einer abstrakten Bewegungsbeschreibung eine reale, koordinierte
Bewegung der Achsen.

Die Übertragung des G-Codes an den Arduino Uno erfolgt über den
Universal G-Code Sender (UGS). UGS verbindet sich seriell über USB mit
dem Controller, zeigt den geladenen Werkzeugpfad als 2D/3D-Vorschau im
Visualizer-Fenster an und übernimmt im sogenannten Batch-Mode die
geordnete, gepufferte Übertragung der einzelnen G-Code-Zeilen an GRBL.

\section{Einrichten und Referenzieren der Maschine}

Bevor ein Werkstück geschnitten werden kann, muss die Maschine in einen
definierten Ausgangszustand gebracht werden. Dazu fahren die Achsen
nacheinander in Richtung der Endschalter, bis diese ausgelöst werden.
Auf diese Weise erkennt die Steuerung die jeweilige Nullposition der
Achsen, unabhängig davon, an welcher Stelle sich der Schneidrahmen zuvor
befunden hat. Diese sogenannte Referenzfahrt (Homing) wird durch den
Befehl \$H ausgelöst und stellt sicher, dass die im G-Code hinterlegten
Koordinaten korrekt auf die tatsächliche Position der Maschine
übertragen werden.

Der aktuelle Zustand der Maschine -- Status (Idle, Run, Hold, Alarm),
Achspositionen sowie Vorschub -- wird während des gesamten Vorgangs im
Digital Read-Out (DRO) von UGS in Echtzeit angezeigt. Vor dem
eigentlichen Schneidvorgang wird zusätzlich der Werkstück-Nullpunkt an
der Maschine festgelegt und in UGS über die Funktion Zero All
zurückgesetzt, sodass die Koordinaten des G-Codes exakt auf die Position
des Schaumstoffblocks abgestimmt sind.

\section{Der Schneidvorgang}

Sobald die Maschine referenziert, der Werkstück-Nullpunkt gesetzt und
der zugehörige G-Code in UGS geladen wurde, kann der eigentliche
Schneidvorgang über die Send-File-Funktion gestartet werden. Die Achsen
bewegen den Heizdraht entsprechend der programmierten Bahn durch den
Schaumstoffblock. Damit der Draht das Material sauber durchtrennt, muss
er auf eine geeignete Temperatur aufgeheizt sein, während die Achsen ihn
mit einer gleichmäßigen, auf das Material abgestimmten Geschwindigkeit
führen.

Die Temperatur des Heizdrahtes wird über das MOSFET-PWM-Schaltmodul
gesteuert. Über ein pulsweitenmoduliertes Signal wird die dem Draht
zugeführte elektrische Leistung und damit seine Temperatur reguliert.
Eine zu niedrige Temperatur führt zu einem unsauberen Schnitt oder zu
erhöhtem Widerstand beim Vorschub, eine zu hohe Temperatur kann das
Material unnötig stark schmelzen und die Schnittkante beschädigen. Die
Abstimmung von Drahttemperatur und Vorschubgeschwindigkeit der Achsen
ist daher entscheidend für die Qualität des Schnittergebnisses.

Während des Schneidvorgangs kann der Fortschritt im Statusbalken sowie
im Konsolenfenster von UGS verfolgt werden. Bei einer Fehlfunktion, etwa
dem unbeabsichtigten Auslösen eines Endschalters, meldet GRBL einen
entsprechenden Alarmzustand (z. B. ALARM:1 bei einem Hard-Limit), der
den Vorgang automatisch anhält. Der Druckknopf dient als einfaches,
lokales Bedienelement, über das ein laufender Vorgang bei Bedarf
zusätzlich manuell unterbrochen werden kann.

\section{Fertigstellung des Werkstücks}

Nach Abschluss der programmierten Schnittbahn fahren die Achsen in der
Regel in eine definierte Ruheposition zurück. Das fertig geschnittene
Werkstück kann anschließend aus dem verbleibenden Schaumstoffblock
entnommen werden. Je nach Anwendung kann eine Nachbearbeitung sinnvoll
sein, etwa das Entfernen kleiner Schmelzgrate an der Schnittkante.

Mit der Entnahme des Werkstücks schließt sich der Kreis vom anfänglichen
Entwurf der Kontur bis zum fertigen, physischen Bauteil. Der
beschriebene Ablauf verdeutlicht, wie die in den vorangegangenen
Kapiteln einzeln vorgestellten Komponenten -- Arduino Uno mit GRBL, CNC
Shield V3, Endschalter, Druckknopf, MOSFET-PWM-Schaltmodul und der
Bediensoftware UGS -- im Zusammenspiel den vollständigen Funktionsumfang
des CNC-Hot-Wire-Foam-Cutters ergeben.

